\section{Introduction}\label{sec:intro}

Alice holds coins at an address she created years ago. She wants them to move when a condition on the chain holds: a height is reached, a balance arrives or the fee market allows it. She does not want to be online when that happens. She is willing to make one move now, which is to post one message. After that a committee of $n$ nodes, of which at most $f$ are Byzantine, must build the transaction from the state of the chain at that later time, sign it under the key of her address and do so once. We call this the \emph{One-Move Spend} problem.

The problem has a trust question at its center. Whatever Alice hands to the committee in her one move, a large enough coalition of nodes holds it jointly and can copy it. What can such a coalition do with the copy, and what stops it? Two published answers exist, and neither solves the problem.

In the first, the committee holds the whole key as a Shamir sharing, and the owner posts one authenticated request~\cite{GS22,TroutPP,HOT}. This is efficient, and a one-message, publicly verifiable dealing can move an existing key into the committee~\cite{cgVSS,Groth21}. Yet a threshold of nodes can then sign any transaction for all time, the limit to one spend is a policy of honest nodes, and the committee must first run a presignature protocol among its members. In the second, the owner keeps one share of the key and posts one partial signature on one message, which the network completes later~\cite{FMMSS25,IkaDocs}. Here the limit to one message is cryptographic. Yet the owner's message consumes a presignature that the network made first, so the owner reads network state before the move, importing an existing key takes two owner transactions and a key registration, and the owner or a contract still builds the transaction. In both forms the signing protocol signs bytes that a caller supplies. Deployed bridges do build transactions inside a replicated state machine, for keys that the committee itself generated~\cite{ckETH,ckBTC,THORChain}.

\paragraph{The idea.} A signature is one point on a line. A Schnorr signature is the value at the challenge $c$ of the line $s(c) = k + c\,x$, and an ECDSA signature is the value at the sighash $h$ of the line $s(h) = h\,k^{-1} + r\,k^{-1}x$. The owner knows both coefficients of the line. The owner therefore deals the line, and does not deal the key. A committee member that holds a Shamir share of each coefficient computes its share of the signature with two field multiplications, and anyone can check that share against public commitments.

The owner's one message is a Shamir dealing with Feldman-style coefficient commitments on the signature curve, with every share encrypted to its node. It comes in two modes.

\begin{itemize}[leftmargin=1.5em]
\item A \emph{sealed point} is the line already evaluated at the sighash of one fully specified transaction. The constant-term commitment of the dealing is not sent. Every verifier computes it from the signature verification equation, as $R + cX$ for Schnorr and as $hG + rX$ in base $R$ for ECDSA, so the committed secret is a valid signature of the address key on that transaction and on nothing else. The owner seals a fee grid of transactions that all spend one pinned input or one account nonce, and the committee chooses one from chain state and releases it.
\item An \emph{open line} is both coefficients. For ECDSA the owner deals $k^{-1}$ and $k^{-1}x$ in base $R = kG$, and the constant terms are forced to $G$ and to the address key $X$. The committee builds the transaction from chain state, fixes its digest with a certificate of a quorum of votes, and only then does each node post one linear partial signature.
\end{itemize}

The two modes give two different answers to the trust question, and \cref{sec:impossibility} shows why open mode cannot confine signatures the way sealed mode does. If an honest committee would sign $m_0$ under one future chain state and $m_1 \neq m_0$ under another, then a coalition that can finish the protocol copies its state after the owner's message and runs the rest of the protocol twice. Erasure does not bind such a coalition. The strongest limit against it is therefore ``only messages that the owner's move already determines'', and two such signatures are harmless only if the chain rejects one of them. Sealed mode reaches this bound: the dealt objects are conflicting signatures, and the chain serializes them. Open mode is on the other side of it: the committee signs a message that later chain state selects, so a qualified coalition can obtain two signatures. The theorem stops there, because two signatures on transactions that spend one pin still give one executed spend. Open mode's one-spend limit is a property of an honest threshold for the reason that follows. A qualified open coalition also holds more than two signatures. The dealt line is the key, and $t$ shares reconstruct it.

\paragraph{Contributions.} We define the One-Move Spend problem and the ideal functionality $\FOMS$ (\cref{sec:problem}). We give the line-dealing protocol in both modes for ECDSA on secp256k1, BIP340 Schnorr and Ed25519. The committee rules make it safe: a request identity that excludes every malleable byte, an anchor rule that fixes one finalized chain state from height reports by the rank rule of oracle aggregation, with retries that consume no nonce, a lock certificate that precedes every partial signature, durable consumption of nonce slots, deterministic transaction construction, and a freeze rule for a committee change (\cref{sec:construction}). We prove the security claims that we make and state the status of each one (\cref{sec:security}). The results are sealed mode under the unforgeability of the signature scheme, open ECDSA by a reduction to Theorem~3 of Groth and Shoup~\cite{GS21}, open Schnorr under the algebraic one-more discrete logarithm assumption in a single-signer game that we define, and the share encryption under interactive CDH after the proof of Groth and Shoup~\cite{GS22}. We prove the impossibility result above (\cref{sec:impossibility}). We report a Lean~4 development of 103 theorems, its port to Natural Lean, a Quint model of the board rules (\cref{sec:formal}), and a Rust prototype with measurements for committees of 9, 16, 64 and 256 nodes (\cref{sec:eval}).

\paragraph{What is not new.} The algebra of the two dealings is old, and we claim none of it. Dealing a finished signature with its commitment fixed by the verification equation is the verifiable signature sharing of Franklin and Reiter~\cite{FR95}, including the ElGamal form in base $r$ that is the ECDSA sealed point. Dealing shares of $k$ and $kx$ for each future signature to signers who answer with one linear partial is the trusted-center scheme of Langford~\cite{Lan95}. The public check of a linear ECDSA partial in base $R$ is in the signing phase of Gennaro and Goldfeder~\cite{GG20}. Fee ladders of pre-signed conflicting transactions are Bitcoin practice~\cite{SHMB20a,SHMB20b,Todd23}. Taking the median of $2f+1$ sorted reports, so that $f$ Byzantine reports cannot move the result outside the honest range, is the aggregation rule of Chainlink off-chain reporting~\cite{OCR21}. What we add is the composition and its analysis. The dealer is the owner of an existing address, bound to the address by a public check. Sealed fee variants are held by a threshold committee and selected from finalized chain state. Open lines sit under commitments that force their constant terms, and no partial precedes the lock. The functionality $\FOMS$ and the impossibility result separate the two modes. \Cref{sec:related} gives the comparison source by source.

\paragraph{Organization.} \Cref{sec:setting} fixes the setting. \Cref{sec:problem} states the problem and $\FOMS$. \Cref{sec:construction} gives the protocol. \Cref{sec:security} states the theorems, their assumptions and their status, and lists what is open. \Cref{sec:impossibility} proves the impossibility result. \Cref{sec:formal,sec:eval} report the formal verification and the measurements. \Cref{sec:related} covers related work and \cref{sec:limits} the limitations.
